\documentclass[14pt,a4paper]{extarticle}
\usepackage[left=3cm,right=2cm,top=2cm,bottom=2cm,headheight=18pt]{geometry}
\usepackage{fontspec}
\setmainfont{Liberation Serif}
\setsansfont{Liberation Serif}
\setmonofont{Liberation Mono}
\usepackage{unicode-math}
\setmathfont{texgyretermes-math.otf}
\usepackage[english]{babel}
\usepackage{amsmath,graphicx,booktabs,longtable,array,enumitem}
\usepackage[normalem]{ulem}
\usepackage{titlesec,fancyhdr}
\usepackage[hidelinks,unicode]{hyperref}
\usepackage{url}
\urlstyle{same}
\setlength{\parindent}{1.27cm}
\setlength{\parskip}{0pt}
\linespread{1}
\setlength{\emergencystretch}{3em}
\clubpenalty=10000
\widowpenalty=10000
\titleformat{\section}{\normalfont\bfseries}{\thesection.}{0.5em}{}
\titlespacing*{\section}{0pt}{1.1em}{0.5em}
\setlist{nosep,leftmargin=1.27cm}
\pagestyle{fancy}
\fancyhf{}
\renewcommand{\headrulewidth}{0pt}
\renewcommand{\footrulewidth}{0pt}
\fancyhead[C]{\ifnum\value{page}>1\relax\thepage\fi}
% Information-analytical material: O‘RQ-682 appended methodology §83 (11–14 pt), contents §84, comparison table §85.
\renewcommand{\normalsize}{\fontsize{14bp}{16.8bp}\selectfont}
\normalsize
\newcommand{\tsize}{\fontsize{12bp}{14.4bp}\selectfont}
\newcolumntype{L}[1]{>{\raggedright\arraybackslash}p{#1}}
\setlength{\LTcapwidth}{\linewidth}
\begin{document}
\noindent English companion translation; the Russian edition is primary.
\begin{flushright}
To the Director of the National Agency\\
for Social Protection under the President\\
of the Republic of Uzbekistan\\[2mm]
Copy: Ministry of Economy and Finance\\
of the Republic of Uzbekistan\\[4mm]
From Abduxoliq Akmal ugli Ashuraliyev,\\
independent researcher\\
(ORCID 0009-0003-5482-5526)\\
Residence: 100017, Tashkent, Yunusobod district,\\
Kashgar massiv, house 16\\
\href{mailto:abdulxoliqashuraliyev@gmail.com}{abdulxoliqashuraliyev@gmail.com}\\[2mm]
7 October 2026
\end{flushright}

\begin{center}\textbf{PROPOSAL}\par
\textbf{to the draft normative legal act being prepared to implement paragraph 14 of Decree of the President of the Republic of Uzbekistan No. PF-193 of 11 September 2026}\end{center}

Under Article 40 of the Constitution of the Republic of Uzbekistan, Article 20 of the Law “On Normative Legal Acts” and Articles 3 and 5 of the Law “On Appeals of Individuals and Legal Entities”, I submit proposals for inclusion in the draft act on the new system for including families in the Social Registry through multidimensional assessment, which the Agency, jointly with the Ministry of Economy and Finance, submits to the Presidential Administration by 1 December 2026.

From 1 January 2027 a family’s category will be determined by a point-based system; under the current Regulation the final assessment and the family’s category are formed by the information module automatically. The proposals ensure that each family’s assessment result follows only from approved rules, is retained with its grounds and is explained to the applicant as Article 24 of the Law “On Personal Data” and Articles 53–54 of the Law “On Administrative Procedures” require, and they remove inconsistencies in the current Regulation on the Procedure for Maintaining the Social Registry that have been demonstrated by calculation. The proposed norms create no new bodies or information systems; implementation within existing funds is proposed subject to the Agency’s assessment (Material 2, Section 10). No personal data of families are requested.

I request consideration and an answer on the merits on each of the following questions:
\begin{enumerate}[leftmargin=1.27cm,label=\arabic*.]
\item Include in the draft the provisions of Section I of Material 1 (points 1–23) or norms securing the same results.
\item Remove the case in which a family with income \textit{equal to} minimum consumption expenditure is recognized under point 32 of the Regulation but assigned to no category of point 39, by amending subpoint “a” of point 33 (Material 1, Section II, point 3).
\item Exclude a negative household-plot area in the formula of point 27 of the Regulation (Section II, point 2).
\item Exclude a fall in counted income when the actual remittance increases, and state the comparison period, in point 26 of the Regulation (Section II, point 1), choosing the option after assessing the aggregate impact of the two options set out in Material 2, Section 10.
\item Ensure that the applicant receives the grounds of a decision on inclusion, refusal, change of category or removal, the objection procedure, and the time limit and procedure for appeal in accordance with Article 24 of the Law “On Personal Data” and Articles 53–54 of the Law “On Administrative Procedures” (Section II, points 4 and 5). If the Agency considers that this obligation is already fulfilled, please state the norm and the means by which the applicant receives this information.
\item Align points 34–35 and subpoint “v” of point 39 of the Regulation with subparagraph “v” of paragraph 3 of Decree PF-258 and item (i) of note (b) to its Annex 1: the Regulation includes in the “family at the poverty boundary” category hardship families whose reduced income is below minimum consumption expenditure, which the Decree does not provide for (Material 2, Section 7).
\item Take this proposal into account in preparing the draft under Article 23 of the Law “On Normative Legal Acts” and inform me when the draft is placed on the portal for discussion of draft normative legal acts under Article 24 of that Law.
\end{enumerate}

Under Article 27 of the Law “On Appeals of Individuals and Legal Entities”, an appeal is considered when all questions raised in it have been considered and the answer contains concrete grounds on each question. Since the draft is submitted by 1 December 2026, I request consideration of the proposal within the one-month period established by Article 28 of that Law and, if additional study is required, to notify me within ten days as that Article provides. Please state which provisions are included in the draft and in what wording, or which current norm already resolves the question, or for what reason a provision is not accepted.

All calculations are reproduced by the attached program. The set of control examples and the program are provided free of charge; I am ready to demonstrate them to the responsible staff.

\noindent Attachments:
\begin{enumerate}[leftmargin=1.27cm,label=\arabic*.]
\item Material 1. Proposals for inclusion in the draft normative legal act.
\item Material 2. Explanatory note with scientific and legal justification.
\item Electronic supplement: control-calculation program, tests and results.
\end{enumerate}
The primary edition of the materials is Russian; Uzbek and English editions are attached. The Uzbek edition of Material 1 gives the amendments in the state language.

\bigskip
\noindent Abduxoliq Akmal ugli Ashuraliyev
\end{document}
